%Chargement des paquets
\usepackage{amsmath}
\usepackage{amsfonts}
\usepackage{amssymb}
\usepackage{enumerate}
\usepackage{mathtools}
\usepackage{bbm}
\usepackage{xparse, etoolbox}
\usepackage{enumerate}
\usepackage{enumitem}
\usepackage{mathabx}
\usepackage{babel}[french]

%environment
\usepackage{amsthm}
\usepackage{keytheorems}
\usepackage{hyperref} 

%Interface théorème
\newkeytheorem{lem}[
	name = Lemme
]

\newkeytheorem{prop}[
	name = Proposition
]

\newkeytheorem{thm}[
	name = Théorème
]

\newkeytheorem{coro}[
	name = Corollaire
]

\renewcommand*{\proofname}{Démonstration}

\theoremstyle{definition}
\newtheorem{defn}{Définition}

\theoremstyle{definition}
\newtheorem*{nota}{Notation}

\theoremstyle{definition}
\newtheorem*{limi}{Liminaire}

\theoremstyle{remark}
\newtheorem{rmq}{Remarque}

\theoremstyle{plain}
\newtheorem*{fait}{Fait}


%Ensembles de nombres
\newcommand{\N} {\mathbb{N}}
\newcommand{\Ne}{\N^\ast}
\newcommand{\Z} {\mathbb{Z}}
\newcommand{\Q} {\mathbb{Q}}
\newcommand{\R} {\mathbb{R}}
\newcommand{\Rb}{\overline{\mathbb{R}}}
\newcommand{\Rp}{\R_+}
\newcommand{\Rm}{\R_-}
\newcommand{\K} {\mathbb{K}}
\newcommand{\Ket} {\mathbb{K^{\ast}}}
\newcommand{\Cx}{\mathbb{C}}

%Ensembles, tribus
\newcommand{\A}{\mathbb{A}}
\renewcommand{\P}{\mathcal{P}}
\newcommand{\T}{\mathcal{T}}
\newcommand{\F}{\mathcal{F}}
\newcommand{\C} {\mathcal{C}}
\newcommand{\ouv}{\mathcal{O}}
\newcommand{\B}{\mathcal{B}}
\newcommand{\M}{\mathcal{M}}
\newcommand{\etag}{\mathcal{E}}
\newcommand{\etagOp}{\etag^{+}(\Omega)}
\newcommand{\etagO}{\etag(\Omega)}
%%Lp avant quotientage
\newcommand{\Lic}{\mathcal{L}^1}
\newcommand{\Lpc}{\mathcal{L}^p}
\newcommand{\Linfc}{\mathcal{L}^{\infty}}
\newcommand{\Lifull}{\Lic(\Omega, \B, m)}
\newcommand{\Lpfull}{\Lpc(\Omega, \B, m)}
\newcommand{\Linffull}{\Linfc(\Omega, \B, m)}
%%Lp après quotientage
\newcommand{\Li}{L^1}
\newcommand{\Ld}{L^2}
\newcommand{\Lp}{L^p}
\newcommand{\Linf}{L^{\infty}}
\newcommand{\Listd}{\Li(\Omega, m)} 
\newcommand{\Ldstd}{\Ld(\Omega, m)} 
\newcommand{\Lpstd}{\Lp(\Omega, m)} 

%Quelques opérateurs
\DeclareMathOperator{\codim}{codim}
\DeclareMathOperator{\rg}{rg}
\DeclareMathOperator{\tr}{tr}
\DeclareMathOperator{\vect}{Vect}
\DeclareMathOperator{\ima}{Im}
\DeclareMathOperator{\id}{Id}
\DeclareMathOperator{\sign}{sign}
\DeclareMathOperator{\ext}{ext}
\newcommand{\ind}{\mathbbm{1}}
\newcommand*{\spr}[2]{\left(\,#1\,\middle|\,#2\,\right)}
\newcommand{\interior}[1]{%
  {\kern0pt#1}^{\mathrm{o}}%
}
\DeclareMathOperator{\card}{card}
\DeclareMathOperator{\ppcm}{ppcm}

%Commandes de pure flemme
\newcommand{\veps}{\varepsilon}
\newcommand{\intab}{\int_{a}^{b}}
\newcommand{\intR}{\int_{-\infty}^{\infty}}
\newcommand{\intinf}{\int_{- \infty}^{+ \infty}}
\newcommand{\equi}{\Leftrightarrow}
\newcommand{\Kev}{$\K$-espace vectoriel }
\newcommand{\Kevs}{$\K$-espaces vectoriels }
\newcommand{\Rev}{$\R$-espace vectoriel }
\newcommand{\Revs}{$\R$-espaces vectoriels }
\newcommand{\Cev}{$\Cx$-espace vectoriel }
\newcommand{\Cevs}{$\Cx$-espaces vectoriels }
\newcommand{\evn}{(E, \norm{.})}
\newcommand{\interf}[2]{[#1,#2]}
\newcommand{\limb}{\lim\limits_}
\newcommand{\lebn}{\lambda_n}
\newcommand{\pp}{\text{~presque partout}}
\newcommand{\derivp}[2]{\frac{\partial #1}{\partial #2}}
\newcommand{\famiu}{(u_i)_{i \in I}}
\newcommand{\sev}{\text{~s.e.v.}}
\newcommand{\Mnp}{M_{n,p}(\K)}
\newcommand{\Mn}{M_{n}(\K)}
\newcommand{\tq}{\text{~t.q.~}}
\newcommand{\mq}{~montrer~que~}
\newcommand{\et}{\quad \text{et} \quad}
\newcommand{\ou}{\text{~ou~}}
\newcommand{\Kx}{\K[X]}


%Commandes de gars chelou
\renewcommand{\subset}{\subseteq}

%Commande restriction, ty duckduckgo
\def\restr#1#2{\mathchoice
              {\setbox1\hbox{${\displaystyle #1}_{\scriptstyle #2}$}
              \restraux{#1}{#2}}
              {\setbox1\hbox{${\textstyle #1}_{\scriptstyle #2}$}
              \restraux{#1}{#2}}
              {\setbox1\hbox{${\scriptstyle #1}_{\scriptscriptstyle #2}$}
              \restraux{#1}{#2}}
              {\setbox1\hbox{${\scriptscriptstyle #1}_{\scriptscriptstyle #2}$}
              \restraux{#1}{#2}}}
\def\restraux#1#2{{#1\,\smash{\vrule height .8\ht1 depth .85\dp1}}_{\,#2}} 

%Quelques normes
\DeclarePairedDelimiter\abs{\lvert}{\rvert}
\DeclarePairedDelimiter\labs{\bigl\lvert}{\bigr\rvert}
\DeclarePairedDelimiter\norm{\lVert}{\rVert}
\DeclarePairedDelimiterXPP\normi[1]{}\lVert\rVert{_1}{\ifblank{#1}{\:\cdot\:}{#1}}
\DeclarePairedDelimiterXPP\normd[1]{}\lVert\rVert{_2}{\ifblank{#1}{\:\cdot\:}{#1}}
\DeclarePairedDelimiterXPP\normp[1]{}\lVert\rVert{_p}{\ifblank{#1}{\:\cdot\:}{#1}}
\DeclarePairedDelimiterXPP\normq[1]{}\lVert\rVert{_q}{\ifblank{#1}{\:\cdot\:}{#1}}
\DeclarePairedDelimiterXPP\norminf[1]{}\lVert\rVert{_\infty}{\ifblank{#1}{\:\cdot\:}{#1}}

%Bracket en angle
\DeclarePairedDelimiter\angl{\langle}{\rangle}

%Set, parenthèses
\newcommand{\set}[1]{\{ #1 \}}
\newcommand{\bigset}[1]{\bigl\{ #1 \bigr\}}
\newcommand{\Bigset}[1]{\Bigl\{ #1 \Bigr\}}
\newcommand{\bigpar}[1]{\bigl( #1 \bigr)}
\newcommand{\Bigpar}[1]{\Bigl( #1 \Bigr)}

%Opitmisation des opérateurs de comparaison
\DeclarePairedDelimiter\tio{o (}{)}
\DeclarePairedDelimiter\gro{O (}{)}


\pagestyle{fancy}

\fancyhead[C]{}
\fancyhead[R]{}

\begin{document}

\underline{Source :} 

\underline{Leçons :}
\begin{itemize}
\item 
\end{itemize}

\section{Le poblème}

On dispose de douze pièce dont une que l'on sait fausse. 
Elle est en fait plus lourde ou plus légère que les 11 autres (qui ont le même poids).
On dispose d'une balance classique (deux plateaux) et on veut déterminer la fausse pièce en 3 pesées.

\section{Résolution \og directe \fg{}}

On sépare les 12 pièces en 3 groupes de 4 et on effectue une première pesée. On doit alors séparer deux cas.
\begin{enumerate}

\item Cas 1 : la 1ère pesée est inégale.
Dans ce cas, on numérote de 1 à 4 les pièces potentiellement plus lourdes
et de 5 à 8 les pièces potentiellement plus légères.
Les pièces de 9 à 12 sont des références.
On effectue la deuxième pesée ainsi :
\[ \set{1,2,3,5} \text{ vs } \set{4, 9, 10, 11} . \]
On distingue à nouveau deux cas.
\begin{enumerate}
\item Cas 1.1 : la 2ème pesée est égale.
Alors les pièces 6, 7 ou 8 sont potentiellement plus légères.
Il suffit donc de comparer les pièces 6 et 7.
\item Cas 1.2 : la 2ème pesée penche à gauche.
Alors les pièces 1, 2 ou 3 sont potentiellement plus lourdes.
Il suffit donc de comparer les pièces 1 et 2.
\item Cas 1.3 : la 2ème pesée penche à droite.
Alors soit la pièce 5 est plus légère soit la 4 est plus lourde.
Il suffit donc de les comparer.
\end{enumerate}
\item Cas 2 : la pesée est égale.
On utilise la même numérotation que pour le cas 1, et on effectue la deuxième pesée ainsi :
\[ \set{9,10,11,1} \text{ vs } \set{2, 3, 4, 5} . \]
On distingue à nouveau deux cas.
\begin{enumerate}
\item Cas 2.1 : la pesée est égale.
Dans ce cas la pièce 12 est différente. 
Une dernière pesée contre la 1 permet de déterminer si elle est plus lourde ou plus légère.
\item Cas 2.2 : la pesée penche à gauche.
Alors les pièces 9, 10 ou 11 sont plus lourdes. On compare donc 9 à 10.
\item Cas 2.3 : la pesée penche à droite.
Il s'agit du symétrique du cas 2.2.
\end{enumerate}
\end{enumerate}

\section{Résolution par signature}

L'idée est la suivante : si on effectue 3 pesées de 4 pièces contre 4, 
chaque pièce peut se trouver, pour chaque pesée, dans une de ces trois situations :
\begin{enumerate}
\item sur le plateau de gauche ;
\item sur le plateau de droite ;
\item non pesée.
\end{enumerate}
Si on encode ces 3 situations par 1, 0 et -1 respectivement, 
on peut donner une signature pour chaque pièce qui est un vecteur de $\set{-1,0,1}^3$.
En définissant la somme vectorielle usuelle on peut définir 12 signatures uniques à conditions de respecter les règles suivantes :
\begin{enumerate}[label=(\roman*)]
\item si un vecteur $(a,b,c)$ est utilisé, le vecteur $(-a,-b,-c)$ n'est pas utilisé ;
\item parmi les 12 signatures 4 comportent un 1 (resp. 0 et -1) en première composante (resp. en deuxième et troisième composante).
\end{enumerate}
La règle (i) assure que deux pièces différentes sont distinguables par leurs signatures (pour rappel, on ne sait pas si la pièce incriminée
est plus lourde ou plus légère, ainsi si la balance penche 3 fois à gauche, les signatures (1,1,1) et (-1,-1,-1) ne sont pas distinguables).
La règle (ii) assure qu'à chaque pesée il y a 4 pièces à gauche, 4 à droite et 4 non pesées.
Il reste à démontrer que l'on peut en effet exhiber une tel jeu de 12 signatures.
On propose le jeu suivant :
\begin{center}
\begin{tabular}{ | c || c | c | c | c | c | c | c | c | c | c | c | c |}
\cline{1-13}
$p_1$ & 1 & 0 & -1 & 1 & 0 & -1 & 1 & 0 & -1 & 1 & 0 & -1 \\ \hline
$p_2$ & 1 & -1 & 0 & 1 & -1 & 0 & 1 & -1 & 0 & -1 & 0 & 1 \\ \hline
$p_3$ & 1 & 1 & 1 & 0 & 0 & 0 & -1 & -1 & -1 & -1 & 1 & 0 \\ \hline
\end{tabular}
\end{center}

où $p_i$ désigne la pesée numéro $i$.

Ainsi, si la balance penche trois fois à gauche (resp. à droite), alors c'est la pièce 1 qui est plus lourde (resp. plus légère), 
et ainsi de suite.

\section{Généralisation}

\subsection{Solution maximale}

Maintenant que l'on dispose d'une solution pour 12 pièces et 3 pesées, peut-on généraliser le problème à $n$ pesées et $x=f(n)$ pièces ?
Il faut déjà déterminer la fonction $f$.
On remarque que notre protocole de construction des signatures consiste en : 
\begin{enumerate}
\item donner toutes les chaînes de $\set{-1,0,1}^n$ soit $3^n$ chaînes ;
\item retirer la chaîne identiquement nulle (qui n'apporte pas d'information) ;
\item sélectionner le nombre maximal de chaînes parmi les $3^n - 1$ restantes en ne prenant jamais une chaîne et sa chaîne opposée.
\end{enumerate}
On peut donc choisir \textbf{au plus} $\dfrac{3^n-1}{2}$ signatures. 
Néanmoins, chaque pesée devant être équilibrée, on doit pouvoir regrouper les signatures en 3 groupes égaux : 
les signatures comportant des 1 en pesée 1, celles comportant des 0, et celles comportant des -1.
Ainsi, une solution maximale comporte $\dfrac{3^n-3}{2}$ signatures. 

\subsection{Construction des solutions}

Pour $n=3$ on dispose déjà d'une solution maximale (on pourrait aussi en exhiber une pour $n=2$).
On propose l'algorithme suivant permettant de construire les signatures pour $n=4$ :
\begin{enumerate}
\item on considère les 12 signatures proposées, à chacune d'entre elles on ajoute soit -1, soit 0 soit 1 pour quatrième composante ;
\item on considère la signature (0,0,0) on lui ajoute 1 (resp. -1) pour quatrième composante ;
\item on considère les 2 signatures non identiquement nulles non utilisées pour $n=3$ 
(au sens où ni la signature ni son inverse ne figure dans le tableau exhibé, ces deux signatures sont logiquement inverses l'une de l'autre),
ici ce sont les signatures (1,-1,1) et (-1,1,-1). On ajoute à ces signatures les deux nombres non uitlisés en 2. soit 0 et -1 (resp. 1). 
\end{enumerate}
Par construction, on génère $3 \cdot \dfrac{3^n-3}{2} + 3 = \dfrac{3^{n+1}-3}{2}$ signatures dont aucune n'est inverse d'une autre.
En itérant cet algorithme, on résout le problème pour $n \geq 2$ quelconque.


\end{document}